% =============================================================================
% information.tex
%
% Satu sumber kebenaran (single source of truth) untuk data penulis, skripsi,
% dan institusi. Semua data di sini dipakai di cover, abstrak, lembar
% persetujuan, prakata, dan halaman-halaman lain yang membutuhkan.
%
% Pemuatan: % =============================================================================
% information.tex
%
% Satu sumber kebenaran (single source of truth) untuk data penulis, skripsi,
% dan institusi. Semua data di sini dipakai di cover, abstrak, lembar
% persetujuan, prakata, dan halaman-halaman lain yang membutuhkan.
%
% Pemuatan: % =============================================================================
% information.tex
%
% Satu sumber kebenaran (single source of truth) untuk data penulis, skripsi,
% dan institusi. Semua data di sini dipakai di cover, abstrak, lembar
% persetujuan, prakata, dan halaman-halaman lain yang membutuhkan.
%
% Pemuatan: % =============================================================================
% information.tex
%
% Satu sumber kebenaran (single source of truth) untuk data penulis, skripsi,
% dan institusi. Semua data di sini dipakai di cover, abstrak, lembar
% persetujuan, prakata, dan halaman-halaman lain yang membutuhkan.
%
% Pemuatan: \input{config/information} dari src/main.tex.
% =============================================================================

% --- Data Penulis ---
\newcommand{\NamaPenulis}{Nama Lengkap}
\newcommand{\NamaPenulisInggris}{Full Name}
\newcommand{\NIM}{A000000000}
\newcommand{\TempatLahir}{Kota}
\newcommand{\TanggalLahir}{1 Januari 2000}
\newcommand{\EmailPenulis}{nama@apps.ipb.ac.id}

% --- Data Pembimbing ---
\newcommand{\NamaPembimbingSatu}{Prof. Dr. Nama Pembimbing Satu, M.Sc.}
\newcommand{\NIPPembimbingSatu}{196001011985031001}
\newcommand{\NamaPembimbingDua}{Dr. Nama Pembimbing Dua, M.Si.}
\newcommand{\NIPPembimbingDua}{197001011998021001}

% --- Data Penguji (opsional, untuk halaman persetujuan) ---
\newcommand{\NamaPengujiSatu}{Dr. Nama Penguji Satu, M.Si.}
\newcommand{\NamaPengujiDua}{Dr. Nama Penguji Dua, M.Si.}

% --- Data Skripsi ---
\newcommand{\JudulSkripsi}{%
  Judul Skripsi Maksimum Tiga Baris, Lima Belas Kata Tidak Termasuk
  Kata Depan dan Kata Sambung
}
\newcommand{\JudulSkripsiInggris}{%
  Title of Thesis Maximum Three Lines, Fifteen Words Excluding Articles
  and Conjunctions
}
\newcommand{\KataKunci}{kata kunci satu, kata kunci dua, kata kunci tiga, kata kunci empat, kata kunci lima}
\newcommand{\Keywords}{keyword one, keyword two, keyword three, keyword four, keyword five}

% --- Data Institusi ---
\newcommand{\Departemen}{Departemen Ilmu Komputer}
\newcommand{\ProgramStudi}{Program Studi S-1 Ilmu Komputer}
\newcommand{\ProgramStudiSingkat}{Ilmu Komputer}
\newcommand{\Fakultas}{Fakultas Matematika dan Ilmu Pengetahuan Alam}
\newcommand{\Universitas}{Institut Pertanian Bogor}
\newcommand{\KotaUniversitas}{Bogor}
\newcommand{\Tahun}{2026}

% --- Data Waktu Pelaksanaan (untuk prakata) ---
\newcommand{\BulanMulaiPenelitian}{Januari}
\newcommand{\TahunMulaiPenelitian}{2026}
\newcommand{\BulanSelesaiPenelitian}{Juni}
\newcommand{\TahunSelesaiPenelitian}{2026}
 dari src/main.tex.
% =============================================================================

% --- Data Penulis ---
\newcommand{\NamaPenulis}{Nama Lengkap}
\newcommand{\NamaPenulisInggris}{Full Name}
\newcommand{\NIM}{A000000000}
\newcommand{\TempatLahir}{Kota}
\newcommand{\TanggalLahir}{1 Januari 2000}
\newcommand{\EmailPenulis}{nama@apps.ipb.ac.id}

% --- Data Pembimbing ---
\newcommand{\NamaPembimbingSatu}{Prof. Dr. Nama Pembimbing Satu, M.Sc.}
\newcommand{\NIPPembimbingSatu}{196001011985031001}
\newcommand{\NamaPembimbingDua}{Dr. Nama Pembimbing Dua, M.Si.}
\newcommand{\NIPPembimbingDua}{197001011998021001}

% --- Data Penguji (opsional, untuk halaman persetujuan) ---
\newcommand{\NamaPengujiSatu}{Dr. Nama Penguji Satu, M.Si.}
\newcommand{\NamaPengujiDua}{Dr. Nama Penguji Dua, M.Si.}

% --- Data Skripsi ---
\newcommand{\JudulSkripsi}{%
  Judul Skripsi Maksimum Tiga Baris, Lima Belas Kata Tidak Termasuk
  Kata Depan dan Kata Sambung
}
\newcommand{\JudulSkripsiInggris}{%
  Title of Thesis Maximum Three Lines, Fifteen Words Excluding Articles
  and Conjunctions
}
\newcommand{\KataKunci}{kata kunci satu, kata kunci dua, kata kunci tiga, kata kunci empat, kata kunci lima}
\newcommand{\Keywords}{keyword one, keyword two, keyword three, keyword four, keyword five}

% --- Data Institusi ---
\newcommand{\Departemen}{Departemen Ilmu Komputer}
\newcommand{\ProgramStudi}{Program Studi S-1 Ilmu Komputer}
\newcommand{\ProgramStudiSingkat}{Ilmu Komputer}
\newcommand{\Fakultas}{Fakultas Matematika dan Ilmu Pengetahuan Alam}
\newcommand{\Universitas}{Institut Pertanian Bogor}
\newcommand{\KotaUniversitas}{Bogor}
\newcommand{\Tahun}{2026}

% --- Data Waktu Pelaksanaan (untuk prakata) ---
\newcommand{\BulanMulaiPenelitian}{Januari}
\newcommand{\TahunMulaiPenelitian}{2026}
\newcommand{\BulanSelesaiPenelitian}{Juni}
\newcommand{\TahunSelesaiPenelitian}{2026}
 dari src/main.tex.
% =============================================================================

% --- Data Penulis ---
\newcommand{\NamaPenulis}{Nama Lengkap}
\newcommand{\NamaPenulisInggris}{Full Name}
\newcommand{\NIM}{A000000000}
\newcommand{\TempatLahir}{Kota}
\newcommand{\TanggalLahir}{1 Januari 2000}
\newcommand{\EmailPenulis}{nama@apps.ipb.ac.id}

% --- Data Pembimbing ---
\newcommand{\NamaPembimbingSatu}{Prof. Dr. Nama Pembimbing Satu, M.Sc.}
\newcommand{\NIPPembimbingSatu}{196001011985031001}
\newcommand{\NamaPembimbingDua}{Dr. Nama Pembimbing Dua, M.Si.}
\newcommand{\NIPPembimbingDua}{197001011998021001}

% --- Data Penguji (opsional, untuk halaman persetujuan) ---
\newcommand{\NamaPengujiSatu}{Dr. Nama Penguji Satu, M.Si.}
\newcommand{\NamaPengujiDua}{Dr. Nama Penguji Dua, M.Si.}

% --- Data Skripsi ---
\newcommand{\JudulSkripsi}{%
  Judul Skripsi Maksimum Tiga Baris, Lima Belas Kata Tidak Termasuk
  Kata Depan dan Kata Sambung
}
\newcommand{\JudulSkripsiInggris}{%
  Title of Thesis Maximum Three Lines, Fifteen Words Excluding Articles
  and Conjunctions
}
\newcommand{\KataKunci}{kata kunci satu, kata kunci dua, kata kunci tiga, kata kunci empat, kata kunci lima}
\newcommand{\Keywords}{keyword one, keyword two, keyword three, keyword four, keyword five}

% --- Data Institusi ---
\newcommand{\Departemen}{Departemen Ilmu Komputer}
\newcommand{\ProgramStudi}{Program Studi S-1 Ilmu Komputer}
\newcommand{\ProgramStudiSingkat}{Ilmu Komputer}
\newcommand{\Fakultas}{Fakultas Matematika dan Ilmu Pengetahuan Alam}
\newcommand{\Universitas}{Institut Pertanian Bogor}
\newcommand{\KotaUniversitas}{Bogor}
\newcommand{\Tahun}{2026}

% --- Data Waktu Pelaksanaan (untuk prakata) ---
\newcommand{\BulanMulaiPenelitian}{Januari}
\newcommand{\TahunMulaiPenelitian}{2026}
\newcommand{\BulanSelesaiPenelitian}{Juni}
\newcommand{\TahunSelesaiPenelitian}{2026}
 dari src/main.tex.
% =============================================================================

% --- Data Penulis ---
\newcommand{\NamaPenulis}{Nama Lengkap}
\newcommand{\NamaPenulisInggris}{Full Name}
\newcommand{\NIM}{A000000000}
\newcommand{\TempatLahir}{Kota}
\newcommand{\TanggalLahir}{1 Januari 2000}
\newcommand{\EmailPenulis}{nama@apps.ipb.ac.id}

% --- Data Pembimbing ---
\newcommand{\NamaPembimbingSatu}{Prof. Dr. Nama Pembimbing Satu, M.Sc.}
\newcommand{\NIPPembimbingSatu}{196001011985031001}
\newcommand{\NamaPembimbingDua}{Dr. Nama Pembimbing Dua, M.Si.}
\newcommand{\NIPPembimbingDua}{197001011998021001}

% --- Data Penguji (opsional, untuk halaman persetujuan) ---
\newcommand{\NamaPengujiSatu}{Dr. Nama Penguji Satu, M.Si.}
\newcommand{\NamaPengujiDua}{Dr. Nama Penguji Dua, M.Si.}

% --- Data Skripsi ---
\newcommand{\JudulSkripsi}{%
  Judul Skripsi Maksimum Tiga Baris, Lima Belas Kata Tidak Termasuk
  Kata Depan dan Kata Sambung
}
\newcommand{\JudulSkripsiInggris}{%
  Title of Thesis Maximum Three Lines, Fifteen Words Excluding Articles
  and Conjunctions
}
\newcommand{\KataKunci}{kata kunci satu, kata kunci dua, kata kunci tiga, kata kunci empat, kata kunci lima}
\newcommand{\Keywords}{keyword one, keyword two, keyword three, keyword four, keyword five}

% --- Data Institusi ---
\newcommand{\Departemen}{Departemen Ilmu Komputer}
\newcommand{\ProgramStudi}{Program Studi S-1 Ilmu Komputer}
\newcommand{\ProgramStudiSingkat}{Ilmu Komputer}
\newcommand{\Fakultas}{Fakultas Matematika dan Ilmu Pengetahuan Alam}
\newcommand{\Universitas}{Institut Pertanian Bogor}
\newcommand{\KotaUniversitas}{Bogor}
\newcommand{\Tahun}{2026}

% --- Data Waktu Pelaksanaan (untuk prakata) ---
\newcommand{\BulanMulaiPenelitian}{Januari}
\newcommand{\TahunMulaiPenelitian}{2026}
\newcommand{\BulanSelesaiPenelitian}{Juni}
\newcommand{\TahunSelesaiPenelitian}{2026}
